\medskip\noindent
\textbf{The finite quadratic calculation, including the eta attachment.}
Write $D_2E=(E\wedge E)_{hC_2}$ and $L_2E=\Sigma^{-1}D_2E$,
with the ordinary interchange action.  Put $M_q=M(\ZZ/q,7)$,
where $q$ is even.  Its integral label chains have generators
$a_7,b_8$ and $\partial b=qa$.  Use the free $C_2$ resolution
with the convention
\[
 \partial(e_i xy)=e_{i-1}\bigl(xy+(-1)^i\tau(xy)\bigr)
                   +(-1)^ie_i\partial(xy).
 \tag{D.2.43q1}
\]
Thus $\tau(aa)=-aa$, $\tau(ab)=ba$, and $\tau(bb)=bb$.
The cells relevant to integral homology through degree fifteen are
\[
\begin{array}{c|l}
14&P=e_0aa\\
15&x=e_0ab,\quad y=e_0ba,\quad z=e_1aa\\
16&T=e_0bb,\quad h=e_1ab,\quad h'=e_1ba,\quad U=e_2aa.
\end{array}
\]
Their boundaries are
\[
\begin{aligned}
 \partial x&=-qP,&\partial y&=qP,&\partial z&=2P,\\
 \partial T&=q(x+y),&\partial h&=x-y+qz,
       &\partial h'&=-\partial h,\qquad\partial U=0.
\end{aligned}
\tag{D.2.43q2}
\]
Cancel the unit incidence in $\partial h$, so that $y=x+qz$,
and set $c_q=x+(q/2)z$.  Then $c_q$ is an integral cycle and
$\partial T=2q c_q$.  It follows that
\[
 H_{14}(D_2M_q;\ZZ)=\ZZ/2\{P\},\qquad
 H_{15}(D_2M_q;\ZZ)=\ZZ/(2q)\{c_q\}.
 \tag{D.2.43q3}
\]
This chain computation applies to spectra: extension of scalars
to $H\ZZ$ is symmetric monoidal, so the derived coinvariant
tensor complex computes the integral homology of the
homotopy-orbit square.  It does not discard the homotopy
information in eta attachments; that information is checked next.

The bottom inclusion gives $D_2S^7\to D_2M_q$.  Its first three
cells are $P_{14},z_{15},U_{16}$, with $\partial z=2P$ and
$\partial U=0$.  The extended-power basis and Nishida relation
in \cite[\S1.4, Lemma~1.4.1 and the following displayed relation,
p.~4]{EH5Behrens2012} identify these mod-two classes as
$Q^7a,Q^8a,Q^9a$ and give
\[
 \operatorname{Sq}_*^2U
  =\operatorname{Sq}_*^2Q^9a
  =\binom72 Q^7a=P.
 \tag{D.2.43q4}
\]
The fifteen-skeleton of $D_2S^7$ is $M(\ZZ/2,14)$;
its fifteenth stable homotopy group is generated by the bottom
eta.  Hence the sixteen-cell $U$ attaches by, and kills, $P\eta$.
The same class remains zero after mapping to any of the quadratic
spectra considered below.

To see that this accounts for the entire possible Hurewicz kernel,
let $E$ be a spectrum with cells in dimensions at least fourteen
and $H_{14}E=\ZZ/2\{P\}$.  The fiber of
$E\to H\ZZ\wedge E$ is $E\wedge\tau_{\ge1}\mathbb S$.
Its lowest possible homotopy group is
\[
 \pi_{15}(E\wedge\tau_{\ge1}\mathbb S)
       =H_{14}E\otimes\pi_1\mathbb S=\ZZ/2,
\]
and its map into $\pi_{15}E$ is the bottom eta action.
Thus the degree-fifteen stable Hurewicz map is surjective with
kernel generated at most by $P\eta$.  Equation~\textup{(D.2.43q4)}
kills that kernel here, proving the natural isomorphisms
\[
 \pi_{15}D_2M_q\xrightarrow{\ \cong\ }
       H_{15}(D_2M_q;\ZZ)=\ZZ/(2q)\{c_q\}.
 \tag{D.2.43q5}
\]
Equivalently, stable unit cancellation gives the relevant
Chang block with attaching matrix $(2,\eta_{14})$, together
with a Moore $2q$ block whose bottom dimension is fifteen.
Any eta component on that Moore block's top attachment is
absorbed by $U$; no hidden homotopy extension is being inferred
from integral chains alone.

The specified comparison $c_{2,4}:M_2\to M_4$ acts on label chains
by $a\mapsto2a$, $b\mapsto b$.  Therefore
\[
 x\longmapsto2x,\qquad z\longmapsto4z,
 \qquad c_2=x+z\longmapsto2x+4z=2c_4.
 \tag{D.2.43q6}
\]
Naturality of the Hurewicz isomorphism computes the map on
$\pi_{15}$ as $\ZZ/4\to\ZZ/8$, $1\mapsto2$.

Now put $K_8=M_4\vee S^8_r$.  Distributing the two smash factors
and taking homotopy orbits gives the actual stable splitting
\[
 D_2K_8\simeq D_2M_4\vee(M_4\wedge S^8_r)\vee D_2S^8_r.
\]
The middle summand is $M(\ZZ/4,15)$, with bottom class $u=ar$;
the last summand starts in dimension sixteen.  Thus
\[
 \pi_{15}D_2K_8=\ZZ/8\{c_4\}\oplus\ZZ/4\{u\}.
 \tag{D.2.43q7}
\]
Attach the nine-cell of
$K=K_8\cup_{i\eta_7+2r}e^9_s$ and filter its square by the
number of factors in the new cell.  The one-new-factor quotient
is $K_8\wedge S^9$; the two-new-factor quotient is $D_2S^9$,
which starts in dimension eighteen.  Through dimension sixteen,
the former quotient has just the cell $as$.  Its attachment is
the orbit image of $a\wedge(i\eta_7+2r)$, hence
$P\eta+2u$.  There is no extra factor of two on the eta term:
the off-diagonal cross term is an induced free $C_2$-orbit,
whose homotopy orbits give one copy, as in the preceding wedge
splitting.  By \textup{(D.2.43q4)} this imposes exactly $2u=0$.
All further cells have dimension at least seventeen and do not
affect $\pi_{15}$.  Consequently
\[
 \begin{aligned}
 \pi_{14}L_2M_2&=\ZZ/4\{c_2\},\\
 \pi_{14}L_2K&=\ZZ/8\{c_4\}\oplus\ZZ/2\{u\},
       & (L_2m_0)_*c_2&=2c_4.
 \end{aligned}
 \tag{D.2.43q8}
\]

Finally, the extra free nine-sphere in the complete target model
adds a free sixteen-cell to the quadratic spectrum, with zero
attachment, and higher-dimensional cells.  It has no effect on
\textup{(D.2.43q8)}.  A target cell of dimension at least ten
first contributes a cross term in dimension at least seventeen,
because the bottom target cell has dimension seven.  Homotopy
orbits add only nonnegative cell dimensions.  Thus the complete
target's quadratic spectrum has the same fifteenth homotopy
group and the same specified map.  This truncation statement
concerns the quadratic layer; it does not truncate the domain
of the target James--Hopf homomorphism.
